\documentclass[10pt,a4paper]{article}
\usepackage[top=2.2cm,bottom=2.2cm,left=2.2cm,right=2.2cm]{geometry}
\usepackage{graphicx}
\usepackage{amsmath,amssymb}
\usepackage{booktabs}
\usepackage{multirow}
\usepackage{makecell}
\usepackage{longtable}
\usepackage{array}
\usepackage{tabularx}
\usepackage[hidelinks]{hyperref}
\usepackage{xurl} % lets long URLs break across lines (loads after hyperref)
\setlength{\emergencystretch}{3em} % avoids slightly overfull lines
\usepackage{xcolor}
\usepackage{colortbl}
\usepackage{caption}
\usepackage{subcaption}
\usepackage{tikz}
\usetikzlibrary{arrows.meta,positioning,fit}
\usepackage{setspace}
\usepackage{enumitem}
\usepackage{float}
\usepackage{times}
\usepackage[T1]{fontenc}
\usepackage[utf8]{inputenc}

\definecolor{toprow}{RGB}{198,239,206}
\definecolor{secrow}{RGB}{197,224,243}
\definecolor{trdrow}{RGB}{255,235,156}
\definecolor{hdrblue}{RGB}{46,64,87}
\definecolor{sigred}{RGB}{154,0,0}

\newcolumntype{C}[1]{>{\centering\arraybackslash}p{#1}}
\newcolumntype{L}[1]{>{\raggedright\arraybackslash}p{#1}}

\newcommand{\inst}[1]{\textsuperscript{\normalfont\small#1}}
\newcommand{\authorrunning}[1]{}
\newcommand{\keywords}[1]{\par\smallskip\noindent\small\textbf{Keywords:}\enspace#1}
\makeatletter
\renewcommand\subsubsection{\@startsection{subsubsection}{3}{\z@}%
  {-8\p@ \@plus -4\p@}{-0.5em}{\normalfont\normalsize\bfseries}}
\makeatother

\onehalfspacing

\begin{document}

\title{\textbf{An Integrated Deep Learning and Machine Learning Framework\\
for Plant Disease Detection and Crop Yield Prediction}}

\author{
  Zahra El Ghourani\inst{1} \quad
  Sireen Hammoud\inst{1} \quad
  Tala Mezher\inst{1} \quad
  Seifedine Kadry\inst{1,2}\\[4pt]
  \small\inst{1}Department of Computer Science and Mathematics,
  Lebanese American University, Beirut, Lebanon\\
  \small\texttt{\{zahra.elghourani, sireen.hammoud, tala.mezher\}@lau.edu}\\[2pt]
  \small\inst{2}Department of Applied Data Science,
  Noroff University College, Kristiansand, Norway\\
  \small\texttt{seifedine.kadry@lau.edu.lb}
}
\date{}
\maketitle
\thispagestyle{empty}

\begin{abstract}
Plant diseases threaten food security, and early detection paired with yield
estimation supports agricultural decision-making. We present a three-stage
framework combining CNN-based disease classification, severity-aware disease
localisation and machine-learning yield prediction. In Stage 1, 38 pretrained CNNs from nine families were benchmarked on the 31,053-image Bangladeshi Crops Disease Dataset (15 classes). A post-hoc audit with exact-file and perceptual hashing showed that about 30\% of the files are duplicate copies and that filename-based grouping had not kept all copies within one split, so results are reported on a de-duplicated test subset ($n=1{,}774$) without retraining. On this subset seven models, led by EfficientNetB7 (93.80\%) and ConvNeXtLarge (93.69\%), are statistically indistinguishable (accuracy 93.4 to 93.8\%; Holm-adjusted McNemar tests), and the compact EfficientNetV2S (93.35\%), used for deployment, belongs to this group.
In Stage~2, a severity-aware
Faster RCNN trained on 2,256 Grounding DINO pseudo-labelled images reached
mAP@0.5 of 49.3\% on a held-out set; against 280 manual annotations, the
pseudo-labels had a mean IoU of 0.389, with systematic single-box-per-image
under-counting. In Stage~3, the yield records are not paired with the images and
the training severity is a fixed per-crop constant, so this stage is a conceptual
demonstration; with countries held out, XGBoost reaches R\textsuperscript{2} =
0.64, no better than a crop-mean baseline (0.64), and severity adds no signal
beyond crop identity. Grad-CAM heatmaps confirm that the classifier attends to
disease-relevant leaf regions.

\keywords{Plant Disease Detection \and Faster RCNN \and Transfer
Learning \and Crop Yield Prediction \and XGBoost \and Precision
Agriculture \and ConvNeXt \and EfficientNet \and Grounding DINO
\and Grad-CAM \and SHAP}
\end{abstract}
